
\begin{table}[H]
\centering
\caption{Heterogeneity: Rural Status and Poverty}
\label{tab:hetero}
\begin{threeparttable}
\begin{tabular}{lcc}
\toprule
& (1) & (2) \\
& Rural/Urban & High/Low Poverty \\
\midrule
ConvShare$^{\text{pre}} \times$ Post & 0.5541$^{***}$ & 0.2221$^{**}$ \\
& (0.0737) & (0.1015) \\[4pt]
ConvShare$^{\text{pre}} \times$ Post $\times$ Rural & -0.5648$^{***}$ & \\
& (0.0942) & \\[4pt]
ConvShare$^{\text{pre}} \times$ Post $\times$ High Poverty & & 0.0122 \\
& & (0.0883) \\[4pt]
County FE & Yes & Yes \\
Year FE & Yes & Yes \\
Observations & 34,435 & 34,156 \\
Clusters (states) & 56 & 56 \\
\bottomrule
\end{tabular}
\begin{tablenotes}[flushleft]
\small
\item \textit{Notes:} Dependent variable: log$(1 + \text{convenience stores})$. Standard errors clustered at the state level in parentheses. * $p<0.10$, ** $p<0.05$, *** $p<0.01$. Rural is defined as USDA Rural--Urban Continuum Code $\geq 4$. High poverty is defined as county poverty rate above the sample median (ACS 2015).
\end{tablenotes}
\end{threeparttable}
\end{table}

